Fix \(P\in\mathcal M_s\).  The nonadaptive design used here is the admissible sequential design
\[
 \pi_t(\cdot\mid\mathcal F_{t-1})=\operatorname{Unif}(\mathcal Z_B),
 \qquad t=1,\ldots,m, \tag{19}
\]
implemented with independent fresh pool randomizations.  Thus its stacked assignment matrix has law \(\Pi_{m,B}\).  For row \(i\), put \(T=T_i=\{i\}\cup S_i\), \(k=|T|\le s+1\), and define \(w\in\mathbb R^T\) by \(w_i=0\) and \(w_j=\operatorname{sign}(\Gamma_{ij})\) for \(j\in S_i\).  The partially penalized Lasso minimizes
\[
 L_i(b)=\frac1{2m}\lVert H_mY_i-Xb\rVert_2^2
       +\lambda_{\mathrm L}\lVert b_{-i}\rVert_1,
 \qquad \lambda_{\mathrm L}=c_{\mathrm L}q\delta, \tag{20}
\]
where the positive universal constant \(c_{\mathrm L}\) will be chosen below.  A minimizer exists.  Indeed, along a minimizing sequence the penalty bounds \(\lVert b_{-i}\rVert_1\); if \(Xe_i\ne0\), the squared-loss bound then also bounds \(b_i\), while if \(Xe_i=0\), coordinate \(b_i\) may be set to zero without changing the objective.  A convergent subsequence and continuity give a minimizer.  The minimizer set is nonempty, closed, and convex, so it has a unique minimum-Euclidean-norm point.  One direct measurability argument is to add \(r^{-1}\lVert b\rVert_2^2\) to (20): the resulting strongly convex coercive objective has a unique minimizer that is continuous in the finite-dimensional data, and as \(r\to\infty\) these minimizers converge to the minimum-norm minimizer of (20).  Indeed, comparison with any minimum-norm minimizer bounds their norms, every cluster point minimizes (20), and the norm comparison selects the unique minimum-norm point.  The selected minimizer is therefore a pointwise limit of Borel maps, so the resulting terminal graph decoder is \(\mathcal F_m\)-measurable.

We first derive the exact-budget population algebra.  By the exact-budget covariance lemma,
\[
 \Sigma:=\operatorname{Cov}(Z_{t\cdot})
 =a\left(I_n-\frac1n\mathbf1_n\mathbf1_n^{\mathsf T}\right),
 \qquad a=p(1-p)\frac n{n-1}. \tag{21}
\]
For any \(T\) of size \(k<n\), direct multiplication verifies
\[
 \Sigma_{TT}^{-1}=a^{-1}\left(I_k+\frac{\mathbf1_k\mathbf1_k^{\mathsf T}}{n-k}\right). \tag{22}
\]
Since \(\Sigma_{jT}=-(a/n)\mathbf1_k^{\mathsf T}\) for \(j\notin T\),
\[
 \left|\Sigma_{jT}\Sigma_{TT}^{-1}w\right|
 =\frac{|\sum_{\ell\in S_i}w_\ell|}{n-k}
 \le\frac{s}{n-s-1}\le\frac16. \tag{23}
\]
The last inequality uses \(s\le n/8\), including \(n=8,s=1\).

Let \(G_T=m^{-1}X_T^{\mathsf T}X_T\).  Apply the balanced-pool uniform-concentration lemma with failure probability \(\eta/8\).  The sample-size hypothesis of the present lemma dominates that lemma's requirement after increasing the universal constant.  Hence, on an event \(E_0\) of probability at least \(1-\eta/8\), simultaneously for all relevant \(T\),
\[
 \lambda_{\min}(G_T)\ge c_0q,
 \qquad \max_{j\in\mathcal U}\lVert Xe_j\rVert_2^2\le2mq. \tag{24}
\]

We next prove rather than assume the support-specific empirical strict-dual bounds.  Fix \(T_0\subsetneq\mathcal U\), \(|T_0|=k_0\le s+1\), \(j\notin T_0\), and \(c\in[-1,1]^{T_0}\), and set
\[
 h=X_{T_0}(X_{T_0}^{\mathsf T}X_{T_0})^{-1}c. \tag{25}
\]
On the corresponding part of (24), this is well defined.  Since \(h\perp\mathbf1_m\), \(h^{\mathsf T}X_j=h^{\mathsf T}Z_j\).  Conditional on the raw columns \(Z_{\cdot T_0}\), independence of the pool rows gives independent Bernoulli coordinates with
\[
 r_t:=\mathbb E(Z_{tj}\mid Z_{tT_0})
 =\frac{B-\sum_{\ell\in T_0}Z_{t\ell}}{n-k_0}.
\]
It follows from \(h^{\mathsf T}Z_{T_0}=c^{\mathsf T}\) that
\[
 \mathbb E(h^{\mathsf T}X_j\mid Z_{\cdot T_0})
 =-\frac{\mathbf1_{k_0}^{\mathsf T}c}{n-k_0}. \tag{26}
\]
The conditional Bernoulli variances are at most \(2q\).  If \(p\le1/2\), then \(r_t\le B/(n-k_0)\le2q\); if \(p>1/2\), apply the same bound to \(1-r_t\) using \(n-B\).  Here \(k_0\le s+1\le n/4\).  Equation (24) and \(\lVert c\rVert_2\le\sqrt{k_0}\) give
\[
 \lVert h\rVert_2^2\le\frac{k_0}{c_0mq},
 \qquad
 \lVert h\rVert_\infty\le\frac{k_0}{c_0mq}. \tag{27}
\]
For the second bound, use \(\lVert X_{tT_0}\rVert_2\le\sqrt{k_0}\) and
\(\lVert(X_{T_0}^{\mathsf T}X_{T_0})^{-1}c\rVert_2\le\sqrt{k_0}/(c_0mq)\).

For clarity, the required conditional Bernstein bound follows elementarily here.  If \(V\) is centered and \(|V|\le M\), expansion of the exponential and \(k!\ge2\,3^{k-2}\) for \(k\ge2\) show
\[
 \mathbb E e^{\lambda V}
 \le\exp\!\left\{\frac{\lambda^2\mathbb E(V^2)}{2(1-|\lambda|M/3)}\right\}
 \quad\text{for }|\lambda|M<3.
\]
Apply this inequality conditionally to the independent variables
\(h_t(Z_{tj}-r_t)\), multiply their moment-generating functions, and optimize in \(\lambda\).  Using (27) and the variance bound gives, for every fixed \(0<u\le1\),
\[
 \Pr\!\left(
  \left|h^{\mathsf T}X_j+\frac{\mathbf1^{\mathsf T}c}{n-k_0}\right|>u
  \middle|Z_{\cdot T_0}\right)
 \le2\exp\!\left(-\frac{c_1mqu^2}{k_0}\right) \tag{28}
\]
on the event in (24), for a universal \(c_1>0\).

Use (28) with \(T_0=T_i\), \(c=w\), \(u=1/6\), and all \(i\) and \(j\notin T_i\).  Equations (23), (26), and a union bound give
\[
 \max_i\max_{j\notin T_i}
 \left|m^{-1}X_j^{\mathsf T}X_{T_i}G_{T_i}^{-1}w\right|
 \le\frac13 \tag{29}
\]
outside an assignment event of probability at most \(\eta/8\).  To bound the active inverse, fix \(\ell\in T_i\), apply (28) with \(T_0=T_i\setminus\{\ell\}\), \(c=w_{T_0}\), and \(u=1/12\), and use the scalar Schur-complement formula.  Its numerator is at most
\[
 |w_\ell|+\frac{|\mathbf1^{\mathsf T}w_{T_0}|}{n-k+1}+\frac1{12}<2,
\]
whereas its denominator is the Schur complement of \(G_{T_0}\) in \(G_T\), which is at least \(\lambda_{\min}(G_T)\ge c_0q\).  A union bound over at most \(n(s+1)\) pairs yields
\[
 \max_i\lVert G_{T_i}^{-1}w\rVert_\infty\le\frac{C_2}{q} \tag{30}
\]
outside another event of probability at most \(\eta/8\).  The requirement
\(mq\ge C(s+1)\log(8n^3/\eta)\) makes both union bounds valid.

We now justify all noise concentration conditional on the complete design.  Under (19), future pool randomizations are independent of the present outcomes and of \(\mathcal F_t\).  Consequently, conditioning additionally on \(Z_{t+1},\ldots,Z_m\) does not change the conditional moment-generating-function bound for \(\varepsilon_{t,i}\).  Backward iteration of the rowwise sub-Gaussian assumption therefore gives, for every vector \(c(Z)\) that is deterministic after conditioning on \(Z\),
\[
 \mathbb E_{P,d}\!\left[\exp\!\left\{\lambda\sum_{t=1}^m c_t(Z)\varepsilon_{t,i}\right\}\middle|Z\right]
 \le\exp\!\left\{\frac{\lambda^2\sigma^2}{2}\sum_{t=1}^mc_t(Z)^2\right\}. \tag{31}
\]
This is the precise point at which nonadaptivity of the upper-bound design is used; no independence across outcome rows is used.

On the assignment event (24), (29), and (30), define the restricted candidate
\[
 \widetilde b_T=\theta_{i,T}
 +G_T^{-1}\left(m^{-1}X_T^{\mathsf T}H_m\varepsilon_i
                   -\lambda_{\mathrm L}w\right),
 \qquad \widetilde b_{T^c}=0. \tag{32}
\]
Because \(X=H_mX\), equation (31) shows that the noise part of coordinate \(\ell\in T\) in (32) is conditionally sub-Gaussian with variance proxy
\[
 \frac{\sigma^2}{m}e_\ell^{\mathsf T}G_T^{-1}e_\ell
 \le\frac{\sigma^2}{c_0mq}. \tag{33}
\]
A union bound over fewer than \(n(s+1)\) coordinates and the term
\(mq\delta^2/\sigma^2\ge C\log(8n^2/\eta)\) make their absolute values at most \(\delta/4\), except on an event of probability at most \(\eta/4\); when \(\sigma=0\), this conclusion is deterministic.  Choose \(c_{\mathrm L}\le(8C_2)^{-1}\).  Equations (30)--(33) then imply, for every \(\ell\in S_i\),
\[
 |\widetilde b_\ell-\theta_{i,\ell}|
 \le\frac\delta4+\lambda_{\mathrm L}\frac{C_2}{q}
 \le\frac{3\delta}{8}<\delta. \tag{34}
\]
By beta-min, every such coordinate has the same nonzero sign as \(\Gamma_{i\ell}\).

Let \(P_T=X_T(X_T^{\mathsf T}X_T)^{-1}X_T^{\mathsf T}\).  For \(j\notin T\), the inactive KKT score at (32) is
\[
 m^{-1}X_j^{\mathsf T}(I-P_T)H_m\varepsilon_i
 +\lambda_{\mathrm L}m^{-1}X_j^{\mathsf T}X_TG_T^{-1}w. \tag{35}
\]
By (24) and (31), the first term is conditionally sub-Gaussian with variance proxy at most \(2\sigma^2q/m\).  The second term has absolute value at most \(\lambda_{\mathrm L}/3\) by (29).  The noise-to-signal sample-size term and a union bound over fewer than \(n^2\) scores make the first term at most \(\lambda_{\mathrm L}/3\) simultaneously, outside an event of probability at most \(\eta/4\).  Hence all inactive dual coordinates have absolute value strictly below one.

On \(T\), equation (32) is exactly the KKT equation for (20) with subgradient \(w\); its coordinate \(i\) is zero because that coordinate is unpenalized.  Equation (34) verifies the active signs, while (35) verifies strict dual feasibility.  Thus convexity proves that \(\widetilde b\) is a minimizer.  It is the unique minimizer: any two minimizers have the same fitted value by strict convexity of the squared loss in \(Xb\); strict dual feasibility forces every other minimizer to vanish on \(T^c\); and positive definiteness of \(G_T\) then forces equality on \(T\).

The three assignment failure probabilities and two noise failure probabilities sum to at most \(7\eta/8<\eta\).  Therefore all rows are recovered simultaneously, so the decoder obtained from the minimum-norm minimizers satisfies \(\widehat A=A\).  The argument is uniform over \(P\in\mathcal M_s\), and the union bounds never use contemporaneous independence across outcome rows.  This proves the displayed error guarantee.  The proof selects exact minimizers only and establishes no finite-precision operation or bit-complexity bound, exactly as stated.